\newcommand{\GaussVMagneticFolio}[3]{%
  \fancypagestyle{GaussVPrinted#1}{%
    \fancyhead{}%
    \fancyhead[C]{\small MAGNETISCHE BEOBACHTUNGEN.}%
    \fancyhead[#2]{\small #1}%
    \fancyfoot[R]{\small #3\hspace{.18\headwidth}}%
    \renewcommand{\headrulewidth}{0pt}%
  }%
  \thispagestyle{GaussVPrinted#1}%
}
\SourceLeaf{572}Beobachtungen geführten Rechnung enthält folgende Tafel in Secunden ausgedrückt:
\Needspace{19\baselineskip}
\begin{center}\small
\textit{Mittleres Schwanken der magnetischen Declination während der drei Jahre}\par
1834---1837.\par\smallskip
\begin{tabular}{l|r|r|r||r||r|r|r||r}
&\multicolumn{4}{c||}{8 Uhr Vormittag}&\multicolumn{4}{c}{1 Uhr Nachmittag}\\\cline{2-9}
&I&II&III&Mittel&I&II&III&Mittel\\\hline
April&74&126&205&147&129&101&264&180\\
Mai&192&124&277&207&158&183&210&185\\
Junius&172&171&199&181&95&151&217&162\\
Julius&213&243&287&250&119&184&252&193\\
August&264&253&269&262&175&165&307&225\\
September&162&325&207&241&172&143&161&159\\
October&116&296&216&222&182&202&242&210\\
November&79&205&308&218&170&173&126&158\\
December&132&324&71&206&184&206&154&182\\
Januar&146&274&138&196&174&212&154&181\\
Februar&116&146&164&143&178&183&129&165\\
März&100&109&366&228&127&153&246&183\\\hline
Mittel&157&229&238&211&156&174&213&183
\end{tabular}
\end{center}

Von den einzelnen Beobachtungen mögen hier noch die \textit{grössten} Schwan\-kungen angeführt werden, die im Laufe der drei Jahre bei den vormittägigen und nachmittägigen Declinationen vorgekommen sind. Jene war am 8. October 1835 um $20'1''$ grösser als am 7. October, und die nachmittägige Declination am 24. April 1836 um $13'0''$ grösser als am vorhergehenden Tage. Dagegen ist auch völlige Gleichheit der vormittägigen oder der nachmittägigen Declination an zweien auf einander folgenden Tagen öfters vorgekommen. In den monatlichen Mittelschwankungen rücken natürlich diese Extreme viel näher zusammen; gleichwohl bleibt die grosse Ungleichheit der einzelnen Monate in dieser Beziehung sehr bemerkenswerth, da nach obiger Übersicht das mittlere Schwanken bei der Vormittagsdeclination im März 1837 die Grösse von $6'6''$ hatte, im December 1836 hingegen nur $1'11''$ betrug.

Ob im Allgemeinen zu einer Tageszeit grössere Schwankungen vorherrschen als zu einer andern, ist aus den Resultaten für unsere beiden Stunden mit Sicherheit noch nicht zu entscheiden. Im Mittelwerth stehen im ersten Jahre beide
\SourceLeaf{573}nahe gleich, in den beiden andern überwiegen die Vormittagsschwankungen, aber der Unterschied der Endresultate aus allen drei Jahren $3'31''$ und $3'3''$ ist zu klein, als dass man ihn durch so wenige Jahre für festgestellt halten dürfte, wiewohl in den Mittelzahlen für die einzelnen Monate in der vierten und achten Columne zehn Monate eine Differenz in demselben Sinn gegeben haben.

Wirft man Vormittags- und Nachmittagsbeobachtungen zusammen, so erhält man folgende mittlere Schwankungen:
\begin{center}\small
\begin{tabular}{l|r|r|r||r}
&Jahr I&Jahr II&Jahr III&Mittel\\\hline
April&108&114&237&164\\
Mai&176&156&245&196\\
Junius&139&161&208&172\\
Julius&173&215&270&223\\
August&224&214&289&244\\
September&167&251&185&204\\
October&152&254&229&216\\
November&133&190&235&191\\
December&160&271&120&195\\
Januar&160&245&146&189\\
Februar&150&166&148&155\\
März&114&133&312&206
\end{tabular}
\end{center}
\Needspace{7\baselineskip}
\begin{center}\small
\textit{Mittelwerthe.}\par\smallskip
\begin{tabular}{l|r|r|r||r}
&Jahr I&Jahr II&Jahr III&Mittel\\\hline
Julius---December&170&234&228&213\\
Übrige Monate&143&167&223&181\\
Ganzes Jahr&158&204&226&198
\end{tabular}
\end{center}

Nach den Zahlen der vierten Columne herrschen in den Monaten Julius---December etwas grössere Schwankungen vor, als in den sechs übrigen, aber die Mittelwerthe $3'33''$ und $3'1''$ sind doch wohl zu wenig verschieden, um daraus mit Sicherheit schliessen zu können, dass jene Jahreszeit grössere Schwankungen mehr begünstigt, zumal da der Unterschied nur hauptsächlich in dem einen Jahre 1835---1836 auf diese Art stark hervorgetreten ist.

Sehr kenntlich ist hingegen die Ungleichheit der Veränderlichkeit in den einzelnen drei Jahren gegen einander gehalten; der Mittelwerth für das dritte Jahr ist fast um die Hälfte grösser, als der Mittelwerth für das erste. Das Ge-\par
\clearpage
\SourceLeaf{574}\GaussVMagneticFolio{568}{L}{}
\noindent neralmittel aus sämmtlichen bisherigen Beobachtungen $3'18''$ könnte daher nach längerer Fortsetzung wohl noch erhebliche Abänderung erhalten.

Dies sind die Resultate, die sich aus den bisherigen täglichen Aufzeichnungen der magnetischen Declination ziehen lassen. Es ist sehr zu wünschen, dass ähnliche Arbeiten an mehrern Orten ausgeführt werden, und an einigen ist seit kurzem schon der Anfang damit gemacht. Wenn, wie in Mailand geschieht, die Beobachtungen nicht nach der Ortszeit, sondern genau gleichzeitig mit den hiesigen angestellt werden, so bietet die Vergleichung der einzelnen Tage noch zu andern Combinationen Gelegenheit dar, welche. wenn sie erst eine etwas beträchtliche Zeit umfassen können, von grossem Interesse sein werden. Die Beobachter, welche es auf eine ähnliche Weise halten, d.\ i.\ ihre Aufzeichnungen zu solchen Zeiten machen, welche mit den hiesigen übereinstimmen, werden daher ersucht, die Resultate aller Tage einzeln mitzutheilen, wobei es jedoch zureicht, sie nur nach Scalentheilen anzugeben, so dass die Verwandlung in Bogentheile erspart werden kann, wenn nur zugleich die nöthigen Reductionselemente bemerkt werden.

\Needspace{20\baselineskip}\vspace{6mm}
\pdfbookmark[1]{Erläuterungen zu den Terminszeichnungen und den Beobachtungszahlen.}{erlaeuterungen-termine}
\begin{center}\large\textbf{\textit{Erläuterungen zu den Terminszeichnungen und den Beobachtungszahlen.}}\par
[\textbf{\textit{Im Auszuge.}}]\end{center}
\begin{center}
\begin{minipage}{.96\linewidth}
\centering\fontsize{8.7}{10.5}\selectfont\bfseries
\hrule height .35pt\relax
\smallskip
Resultate aus den Beobachtungen des magnetischen Vereins.\quad 1836.\quad V.\par
\smallskip
\hrule height .35pt\relax
\end{minipage}
\end{center}

Es werden hier ---\quad---\quad---\quad---\quad--- die graphischen Darstellungen der Variationsbeobachtungen von sechs Terminen gegeben, zusammen sechsundvierzig Curven aus vierzehn verschiedenen Beobachtungsörtern: Berlin, Breda, Breslau, Catania, Freiberg, Göttingen, Haag, Leipzig, Mailand, Marburg, Messina, München, Palermo und Upsala.
\par\smallskip\noindent\hspace*{.07\linewidth}\makebox[.93\linewidth][s]{---\hfill---\hfill---\hfill---\hfill---\hfill---\hfill---\hfill---\hfill---\hfill---\hfill---\hfill---\hfill---\hfill---\hfill---}\par\smallskip
Am 28. November 1835 und während der folgenden Nacht wurden die Beobachtungen in Palermo durch einen überaus heftigen Siroccowind sehr gestört, so dass sie einmal sogar auf anderthalb Stunden unterbrochen werden mussten:
\par
\clearpage
\SourceLeaf{575}\GaussVMagneticFolio{569}{R}{72}
\noindent zu vielen Sätzen konnten nur einzelne unzuverlässige Bestimmungen erhalten werden. Es ist daher zu vermuthen, dass viele der sich ergebenden Schwankungen keine reell magnetische Bewegungen gewesen sind. Wir haben indessen doch diese Curve nicht ausschliessen wollen, da der letzte Theil, vom Vormittage des 29. November, wo der Sturm sich ziemlich gelegt hatte, eine ganz befriedigende Harmonie mit den nordlichen Beobachtungsorten zeigt.

Es mag bei dieser Gelegenheit hier noch bemerkt werden. dass nach allen sonstigen Erfahrungen die heftigsten Sturmwinde ohne alle Wirkung auf die Magnetnadel sind, wenn nur durch Dichtigkeit des Locals und Kastens ihr unmittelbarer mechanischer Einfluss hinlänglich abgewehrt ist. Sehr oft ist im Göttinger magnetischen Observatorium während des heftigsten Sturmes von aussen ein äusserst ruhiges Verhalten der Nadel oder ein sehr gleichförmiges Fortschreiten der Variation beobachtet. Wer jedoch nach solchen Erfahrungen gerade umgekehrt vermuthen wollte, dass Stürme in der Atmosphäre den magnetischen Potenzen lähmend entgegenwirkten, würde durch den Hergang des Januartermins 1836 widerlegt werden. Während dieses Termins herrschte in Göttingen und an mehrern andern Beobachtungsorten ein sehr heftiger Sturm, und mehrere auswärtige Beobachter äusserten bei der Einsendung der Resultate die Besorgniss, dass diesmal jenes Umstandes wegen wohl eine geringe Übereinstimmung in den ungemein starken Bewegungen Statt finden werde: gleichwohl war in diesem Termine, wie die Darstellung ---\quad---\quad---\quad--- zeigt, die Harmonie der Curven von den verschiedenen Beobachtungsorten so vollkommen, dass man sie bewundernswürdig nennen müsste, wenn sie nicht nach so vielen Erfahrungen etwas Gewohntes geworden wäre. Eben so wenig wie Stürme haben Gewitter, selbst wenn sie nahe waren, nach mehrern hier und an andern Orten vorgekommenen Erfahrungen, einen erkennbaren Einfluss auf die Magnetnadel gezeigt\expandafter\def\csname HistoricalNoteText575\endcsname{Natürlich ist hier nicht die Rede von dem Falle, wo die atmosphärische Electricität vermittelst eines Zuleitungsdrahts durch einen die Nadel umgebenden Multiplicator zur Erde geführt wird.}\HistoricalNote{575}.

Ein im August 1836 eingelaufenes Schreiben des Hrn.\ \textsc{von Humboldt} enthielt die Nachricht, dass vom 10---18. August zu Reikiavik auf Island die magnetische Variation durch einen geübten französischen Astronomen Hrn.\ \textsc{Lottin} mit einem \textsc{Gambey}schen Apparat ununterbrochen von Viertelstunde zu Viertelstunde beobachtet werden würde, und den Wunsch, dass an einem oder einigen jener Tage correspondirende Beobachtungen mit Magnetometern gemacht werden
\par
\clearpage
\SourceLeaf{576}\GaussVMagneticFolio{570}{L}{}
\noindent möchten. Es wurde dem zu Folge ein ausserordentlicher Termin auf den 17---18. August veranstaltet, und so viel die Kürze der Zeit verstattete, auswärtige Mitglieder unsers Vereins zur Theilnahme eingeladen. Dieser ausserordentliche Termin ist in Upsala, Haag, Göttingen, Berlin, Leipzig und München ganz auf die in den ordentlichen Terminen eingeführte Art abgehalten, und wenn die ---\quad---\quad---\quad---\quad--- graphisch dargestellten Beobachtungen recht interessante Bewegungen zeigen, so müssen wir nur bedauern, dass der ---\quad---\quad---\quad--- für die französischen Isländer Beobachtungen offen gehaltene Platz hat leer bleiben müssen, da wir über den Erfolg dieser französischen Beobachtungen Nichts haben in Erfahrung bringen können.

Der Septembertermin bietet eine Erfahrung dar, die hier etwas ausführlich erwähnt werden mag, da sie das oben [S.\ 552 d.\ B.] bemerkte auf eine lehrreiche Weise bestätigt. Im Protocoll der Marburger Beobachtungen, die dasmal in Abwesenheit des Hrn.\ Prof.\ \textsc{Gerling} ohne dessen persönliche Theilnahme ausgeführt waren, fanden sich für $12^{\mathrm h}5'$ nur ganz unordentlich laufende Zahlen aufgeführt, die gar kein Resultat geben; für $12^{\mathrm h}10'$ erscheint auf einmal eine um $30{,}54$ Scalentheile grössere Zahl, als für $12^{\mathrm h}0'$. ---\quad---\quad---\quad---\quad---\quad---\quad---\quad---\quad---\quad---\quad---\quad--- Diese Erscheinung erregte die Vermuthung, dass um die Zeit $12^{\mathrm h}5'$ eine Spinne die freie Bewegung der Nadel durch Anknüpfung eines Fadens gehemmt habe, und diese Vermuthung erhielt noch eine verstärkte Wahrscheinlichkeit durch den Umstand, dass von $12^{\mathrm h}10'$ bis zu Ende die Bewegungen der Nadel zwar denen, welche die Beobachtungen von andern Orten ergaben, ganz ähnlich, aber verhältnissmässig viel kleiner hervortraten, als man nach den Erfahrungen aus andern Terminen hätte erwarten müssen. Hr.\ Prof.\ \textsc{Gerling} wurde deshalb gebeten, nach seiner Zurückkunft nach Marburg eine genaue Besichtigung des Apparats vorzunehmen, wovon das Resultat aus einem Schreiben des Hrn.\ Prof.\ \textsc{Gerling} vom 12. November hier noch beigefügt werden mag.

Die Untersuchung wurde am 5. November vorgenommen, bis wohin seit dem Septembertermine Niemand wieder in das Beobachtungszimmer gekommen war. Zuerst wurde der Stand der Nadel bestimmt und gefunden
\begin{center}
\begin{tabular}{rrc l}
$3^{\mathrm h}$&$33'$&.\quad.\quad.\quad.\quad.\quad.&$445{,}63$ Scalentheile\\
&35&.\quad.\quad.\quad.\quad.\quad.&$445{,}73$\\
&37&.\quad.\quad.\quad.\quad.\quad.&$445{,}71$
\end{tabular}
\end{center}

\par
\clearpage
\SourceLeaf{577}\GaussVMagneticFolio{571}{R}{72\textsuperscript{*}}
Hierauf wurde die Nadel mit Hülfe des sogen. Beruhigungsstabes in mässige Schwingungen versetzt, und daraus eine Schwingungsdauer von 17 Secunden gefunden, neun Secunden geringer als die sonst bekannte Schwingungsdauer. Als darauf der Deckel des Kastens vorsichtig abgehoben wurde, bemerkte man an dessen unterer Fläche eine sehr kleine lebendige Spinne; auch glaubte man einen daran hängenden, wiewohl kaum bemerkbaren Faden zu gewahren: man fand ferner im Kasten eine Anzahl kleiner schwarzer punktartiger Körper, die sich unter dem Mikroskop als Mückencadaver erwiesen, imgleichen zuletzt in einer Ecke des Kastens ein förmliches unversehrtes Gewebe, von solcher Feinheit, dass es ohne den Wiederschein der Lichter schwerlich erkennbar gewesen wäre. Nach allen Umständen konnte man nur annehmen, dass die Spinne schon seit längerer Zeit ihren Aufenthalt im Kasten gehabt habe.

Nachdem dann noch der Magnetstab auf allen Seiten mit dem Finger umfahren war, ergaben neue Beobachtungen der Schwingungsdauer genau wieder den alten Werth von 26 Secunden. Auch fand sich der Stand wieder bedeutend kleiner als vorher, nemlich
\begin{center}
\begin{tabular}{rrc l}
$4^{\mathrm h}$&$43'$&.\quad.\quad.\quad.\quad.\quad.&$431{,}15$ Scalentheile\\
&45&.\quad.\quad.\quad.\quad.\quad.&$431{,}46$\\
&47&.\quad.\quad.\quad.\quad.\quad.&$431{,}12$
\end{tabular}
\end{center}
Indessen können natürlich diese Standbeobachtungen zu einer genauen Bestimmung, um wieviel die Stellung durch das jetzt weggeschaffte Hinderniss verfälscht gewesen war, nicht dienen, da die etwaige Veränderung der Declination während der mehr als eine Stunde betragenden Zwischenzeit unbekannt blieb.

Bei dieser Gelegenheit mag hier noch ein zweiter Vorfall ähnlicher Art erwähnt werden. Die Schwingungsdauer des Magnetstabes in Breslau, welche im März 1836 beinahe 32 Secunden betrug, hatte von da bis zum November ganz allmählich sich vergrössert, und während dieser Zeit zusammen um etwa $0''4$ zugenommen. Dies ist ganz in der Ordnung, da alle Magnetstäbe (wenn gleich, nach Maassgabe der ungleichen Härtung des Stahls und anderer Umstände, in sehr ungleichen Verhältnissen) im Laufe der Zeit etwas von ihrer Kraft verlieren. Allein vom November 1836 bis Januar 1837\expandafter\def\csname HistoricalNoteText577\endcsname{Vermuthlich waren in der \textit{Zwischenzeit} keine Bestimmungen der Schwingungsdauer gemacht.}\HistoricalNote{577} hatte im Gegentheil wieder eine
\par
\clearpage
\SourceLeaf{578}\GaussVMagneticFolio{572}{L}{}
\noindent Abnahme der Schwingungsdauer von $1''27$ Statt gefunden. und Hr.\ Prof.\ \textsc{von Boguslawski}, welcher mir diesen auffallenden Umstand in einem Schreiben vom 5. März anzeigte, schien geneigt, dies zum Theil auf eine vergrösserte Intensität des Erdmagnetismus zu schieben. Mir jedoch schien nicht zweifelhaft, dass der Grund dieses Phänomens in der nächsten Umgebung des Magnetstabs gesucht werden müsse, wahrscheinlich in nicht ganz freier Beweglichkeit desselben, obwohl von einem Spinnefaden in Gemässheit der Marburger Erfahrungen eher eine bedeutend stärkere Wirkung zu erwarten gewesen sei. Diese Vermuthung fand auch Hr.\ \textsc{von Boguslawski} bestätigt. Unter dem 21. März erwiederte er: `Die Ursache der Änderung der Schwingungsdauer haben Sie richtig errathen. Der Kasten war durch Zufall etwas seitwärts geschoben, so dass der Rand des kleinen Loches, durch welches der Faden von oben eintritt, dem letztern nahe gekommen war, jedoch keineswegs bis zur Berührung. Dennoch müssen einige feine Fasern bis zum Rande gereicht haben, denn seitdem der Faden wieder durch die Mitte des Loches geht, ist auch die Schwingungsdauer wieder nahe dieselbe wie früher.'

Über die Bewegungen selbst, die hier aus sechs Terminen dargestellt werden, mögen einige Bemerkungen hier noch Platz finden.

In den drei Sommerterminen ---\quad---\quad---\quad---\quad---\quad---\quad--- sieht man durch alle grossen Anomalien doch auch die tägliche regelmässige Bewegung durchscheinen, in so fern die Curven in den Nachmittagsstunden aufsteigen, und in den folgenden Vormittagsstunden niedersteigen; in den drei Winterterminen hingegen ---\quad---\quad---\quad---\quad---\quad---\quad--- ist davon kaum noch etwas zu erkennen. Dass das Regelmässige von dem Unregelmässigen überragt wird, oder ganz darin untergeht, ist in der That nach allen unsern Erfahrungen ein sehr gewöhnlicher Hergang: es sind jedoch in den Jahren 1834 und 1835 auch einige Termine vorgekommen, wo der regelmässige Gang durch gar keine grössere Anomalieen verdunkelt wurde, während kleine nie fehlten.

Was aber die anomalen Bewegungen so merkwürdig macht, ist die ausserordentlich grosse, gewöhnlich bis auf die kleinsten Nuancen sich erstreckende Übereinstimmung an verschiedenen Orten, ja meistens an sämmtlichen Orten, nur in ungleichen Grössenverhältnissen. Es würde ganz unnöthig sein, diese Harmonie im Einzelnen nachzuweisen: der Anblick unserer sechs Termindarstellungen spricht hier schon hinlänglich für sich selbst.

\par
\clearpage
\SourceLeaf{579}\GaussVMagneticFolio{573}{R}{}
Für jetzt kann es noch gar nicht unser Beruf sein, diese räthselhafte Hieroglyphenschrift der Natur zu entziffern: wir müssen vorerst unser Bestreben nur sein lassen, Abschriften von dem, was sich darbietet, zu sammeln, und denselben immer mehr Zuverlässigkeit, Treue und Mannigfaltigkeit zu verschaffen: reichem Stoff wird, wie wir zuversichtlich hoffen dürfen, dereinst auch die Entzifferung nicht fehlen. Inzwischen wird es verstattet sein, einige Bemerkungen beizufügen, die zu einer richtigern Beurtheilung beitragen können.

Zuvörderst darf nicht vergessen werden, dass alle solche Anomalieen vergleichungsweise nur geringe Abänderungen oder Zusätze zu der grossen erdmagnetischen Kraft sind (oder genau zu reden, zu dem horizontalen Theile derselben); dass wir zwischen jenen und dieser wohl unterscheiden müssen, und dass, wie die Sache bis jetzt steht, Nichts uns nöthigt, beide gleichen oder gleichartigen Ursachen zuzuschreiben. Immerhin mag man es für wahrscheinlich halten --- was wir ganz auf sich beruhen lassen --- dass jene Anomalieen Wirkungen von electrischen Strömungen oder Ausgleichungen, vielleicht weit ausserhalb der Atmosphäre, sein können: man braucht deshalb doch die ältere Vorstellung noch nicht fahren zu lassen, dass die Hauptkraft in dem festen Theile des Erdkörpers selbst ihren Sitz habe, oder vielmehr die Gesammtwirkung aller magnetisirten Theile des Erdkörpers sei. Wäre, nach der Meinung einiger Naturforscher, das Innere der Erde noch in flüssigem Zustande, so böte die immer fortschreitende Erhärtung und die daraus folgende Verdickung der festen Erdrinde die natürlichste Erklärung der Säcularänderungen der magnetischen Kraft dar.

Wir verlassen jedoch lieber den lockern Boden der Hypothesen, und kehren zu den Thatsachen zurück. Bei weiten die meisten Anomalieen finden wir kleiner an den südlichern Beobachtungsorten, grösser an den nordlichern. So beträgt z.\ B.\ das merkwürdige Aufsteigen am 30. Januar 1836 zwischen $9^{\mathrm h}25'$ und $9^{\mathrm h}40'$, auf Bogentheile reducirt, in Catania 6, in Mailand 12, in München $13\frac12$, in Leipzig 16, in Marburg 20, in Göttingen 26, im Haag 29 Minuten. Von dieser Ungleichheit ist nun zwar etwas abzurechnen wegen des Umstandes, dass an den nordlichern Punkten, wo der horizontale Theil der erdmagnetischen Kraft selbst eine geringere Intensität hat, als an den südlichern, gleiche störende Kräfte eine stärkere Wirkung hervorbringen müssen als an den letztern: allein der Unterschied der Intensitäten vom Haag bis Catania ist im Vergleich mit den beobachteten Ungleichheiten nur gering, und es steht also fest, dass die Energie der
\par
\clearpage
\SourceLeaf{580}\GaussVMagneticFolio{574}{L}{}
\noindent damaligen störenden Kraft desto schwächer war, je weiter nach Süden wir ihre Wirkung verfolgen. Bei aller Unwissenheit, in der wir uns in Beziehung auf das Wesen solcher störenden Kräfte befinden, können wir doch nicht zweifelhaft sein, dass die Quelle einer jeden irgendwo im Raume ihren bestimmten Sitz haben müsse, und so wie wir den Sitz von derjenigen, welche die erwähnten Erscheinungen hervorbrachte, nothwendig nordlich oder nordwestlich von den Beobachtungsorten annehmen müssen (ohne nach so wenigen Datis eine nähere Bestimmung zu wagen), so scheinen überhaupt die nordlichsten Gegenden der Hauptheerd zu sein, von wo die meisten und grössten Wirkungen ausgehen, so weit man nemlich auf Erfahrungen aus einem gegen die ganze Erdfläche doch nur kleinen Umkreise schon derartige Folgerungen stützen darf.

Betrachtet man den bis jetzt vorliegenden Stoff genauer, so finden sich doch bei den verschiedenen Bewegungen, die nach einander vorkommen, rücksichtlich ihrer Grössenverhältnisse an verschiedenen Orten, auch wenn sonst die Ähnlichkeit ganz unverkennbar ist, bedeutende Verschiedenheiten: so ist z.\ B.\ oft von zwei kurz nach einander folgenden Hervorragungen an einem Orte die erste die grössere, an einem andern Orte umgekehrt. Wir werden daher genöthigt, anzunehmen, dass an demselben Tage und in derselben Stunde viele Kräfte zugleich thätig sind, die vielleicht ganz von einander unabhängig sein, sehr verschiedene Sitze haben, und deren Wirkungen an verschiedenen Beobachtungsörtern nach Maassgabe der Lage und Entfernung in sehr ungleichen Verhältnissen sich vermengen, oder, indem die eine zu wirken anfängt, bevor die andere aufgehört hat, in einander eingreifen können. Die Lösung der Verwicklungen, welche dadurch in die Erscheinungen an jedem einzelnen Orte kommen, wird unstreitig sehr schwer sein: gleichwohl dürfen wir zuversichtlich hoffen, dass diese Schwierigkeiten nicht immer unüberwindlich bleiben werden, wenn die Theilnahme an den gleichzeitigen Beobachtungen eine noch viel ausgedehntere Verbreitung erhalten haben wird. Es wird der Triumph der Wissenschaft sein, wenn es dereinst gelingt, das bunte Gewirre der Erscheinungen zu ordnen, die \textit{einzelnen} Kräfte, von denen sie das zusammengesetzte Resultat sind, auseinander zu legen, und einer jeden Sitz und Maass nachzuweisen.

Nicht ganz selten findet man auch bei einzelnen Orten eine kleine Aufwallung, wozu an den übrigen Orten sich kein Gegenstück erkennen lässt. Es würde aber zu gewagt sein, dergleichen sofort für eine blos locale magnetische Einwir-\par
\clearpage
\SourceLeaf{581}\GaussVMagneticFolio{575}{R}{}
\noindent kung zu erklären. Bei einer so grossen Menge von Zahlen kann in der That zuweilen einmal ein Irrthum vorgefallen sein. Öfters sind uns solche Fälle vorgekommen, wo das Nachsehen der Originalbeobachtungen, wenn dieselben in unsern Händen waren, einen Rechnungsfehler in der Reduction oder einen offenbaren Schreibfehler erkennen liess. Zuweilen hat in ähnlichen Fällen, wo wir aber nur den Auszug aus den Beobachtungen zu Händen hatten, die Anzeige eines solchen Verdachts bei dem Einsender einen gleichen Erfolg gehabt. Da jedoch unthunlich ist, alle dergleichen Fälle immer erst durch Briefwechsel zu discutiren, so werden diejenigen Theilnehmer, welche nicht die Originalbeobachtungen selbst einsenden, in Beziehung auf solche Stellen in den ihre Beobachtungen darstellenden Curven, wie z.\ B.\ bei Leipzig am 26. November 1836 für $6^{\mathrm h}15'$ Göttinger Zeit, die Originalbeobachtungen selbst nachzusehen Anlass nehmen können: Berichtigungen, die auf solche Art hervorgehen, sollen dann nachträglich angezeigt werden. Völlige Gewissheit hat man jedoch in Beziehung auf solche Stellen, die nur auf Einem Beobachtungssatze beruhen, auch dann noch nicht, wenn die Originalpapiere keinen Fehler bestimmt nachweisen, da es auch einem nicht ganz ungeübten Beobachter wohl einmal begegnen kann, in demselben Satze wiederholt unrichtige Zehner der Scalentheile niederzuschreiben. Durch eine solche, freilich etwas gewagte Conjectur würde sich z.\ B.\ die oben bemerkte Zahl von $11.69$ auf $6.69$ bringen, also in die übrigen hereintretend machen lassen.

Für local im \textit{engsten} Sinn würde man übrigens eine solche isolirte Aufwallung, auch wo die Thatsache keinem Zweifel mehr unterliegt, immer noch nicht halten dürfen. Wie die Quelle jeder Anomalie irgendwo ihren Sitz haben muss, so kann von dieser oder jener der Sitz auch einmal in der Nähe eines der Beobachtungsörter selbst sein: ist eine solche Kraft an sich nur schwach, so kann ihre Wirkung an dem nächsten Orte, eben wegen der Nähe, augenfällig sein, und verschwindend (d.\ i.\ uns nicht mehr erkennbar) an allen übrigen Orten, \textit{wo beobachtet wird}, eben weil diese schon \textit{zu} entfernt sind. Es scheint daher, bis jetzt wenigstens, gar kein Grund vorhanden zu sein, unter den Anomalieen andere als quantitative Verschiedenheiten anzunehmen. Zugleich aber knüpft sich hieran die Folgerung, dass es in manchen Fällen sehr nützlich sein kann, wenn zwei oder mehrere Beobachtungsörter in nur mässiger Entfernung von einander liegen. Es wäre z.\ B.\ recht erwünscht gewesen, wenn in Augsburg (wo jetzt regelmässig Theil an den Terminsbeobachtungen genommen wird) schon der Septembertermin
